\phantomsection
\addcontentsline{toc}{chapter}{Editor's Preface}
\fancyhead[C]{{\LARGE Editor's Preface}}
\par\noindent
\begin{quoting}
	And Maccabaeus and they that were with him, the Lord leading them on, recovered the temple and the city; and they pulled down the altars that had been built in the marketplace by the aliens, and also the walls of sacred inclosures. And having cleansed the sanctuary they made another altar of sacrifice; and striking stones and taking fire out of them, they offered sacrifices, after they had ceased for two years, and burned incense, and lighted lamps, and set forth the shewbread. And when they had done these things, they fell prostrate and besought the Lord that they might fall no more into such evils; but that, if ever they should sin, they might be chastened by him with forbearance, and not be delivered unto blaspheming and barbarous heathen. ---2 Maccabees 10:1-4
\end{quoting}

\lett{W}{orshiping} the Lord in the beauty of holiness in the Sacrifice of the Mass is the greatest joy of the Christian. Attending to the morning and evening sacrifices of the Daily Office, he finds the apex of his worship, praise, and thanksgiving in the Holy Communion. Because of this, we have taken great care in the formation of this Missal.\\

Based on the 2025 Proposed Book of Common Prayer, it provides everything necessary for the celebration of the Mass for any occasion throughout the year. This also accounts, with great concern and detail, for every peculiarity of the English liturgical tradition, while retaining the great riches gained by the Missal tradition as celebrated in the Antiochian Western Rite Vicariate.\\

In the Mass, we remember that we are sons of the Resurrection, redeemed by the ever-sufficient Blood of Jesus. Because of this, we chose to publish this Missal on 3 May 2026: The Festival of the Invention of the Holy Cross and the Third Sunday after Easter.\\

May this Book provide blessing and great use for every priest who celebrates the Mass and for all laymen who devote themselves to it.\\

\rightline{On the Feast of the Invention of the Holy Cross}

\vspace{1\baselineskip}

\rightline{\emph{Pax Christi,}}

\vspace{2\baselineskip}

\rightline{Mr. Augustine Watson, MSt}

\rightline{\emph{Chief Editor, Apologia Anglicana}}